% answers_short.tex —— 1990 年·全国卷·理科 · 简短答案速查
% 由 tools/make_short_answers.py 自 answers.tex 的 \gksolhead/\gkanswer 提取，
% 并与 1990/1990_数学真题_全国理_答案要点.md 逐题比对后生成，勿手改（改源 MD 后重跑）。

\begin{gknotice}
\textbf{使用说明}：本节只给最终答案便于快速核对，完整推导见解析版。答案为编写者独立推导，\textbf{非官方标准答案}。
\end{gknotice}

\gkpart{选择题}
\begin{center}\begin{tabular}{cccccccccccccccc}
\toprule
题号 & 1 & 2 & 3 & 4 & 5 & 6 & 7 & 8 & 9 & 10 & 11 & 12 & 13 & 14 & 15 \\
\midrule
答案 & A & B & D & C & C & B & A & D & B & D & C & B & B & C & D \\
\bottomrule
\end{tabular}\end{center}

\gkpart{填空题}
\begin{itemize}[leftmargin=2.6em]
  \item[\textbf{16.}] $y=\pm\dfrac{16}{5}$
  \item[\textbf{17.}] $-20$
  \item[\textbf{18.}] $2$
  \item[\textbf{19.}] $\sqrt2+\dfrac12$
  \item[\textbf{20.}] 待核（靠 $A$ 侧:其余=$7:5$；$V_1,V_2$ 对应未标）
\end{itemize}

\gkpart{解答题}
\begin{itemize}[leftmargin=2.6em]
  \item[\textbf{21.}] $(0,4,8,16)$ 或 $(15,9,3,1)$
  \item[\textbf{22.}] $\dfrac{24}{7}$
  \item[\textbf{23.}] $60^\circ$
  \item[\textbf{24.}] $z=\pm(\sqrt{1+a}-1)$；当 $0\le a\le1$ 时另有 $z=\pm i(1\pm\sqrt{1-a})$
  \item[\textbf{25.}] 椭圆方程为 $\dfrac{x^2}{4}+y^2=1$；所求点为 $(\pm\sqrt3,-\dfrac12)$
  \item[\textbf{26.}] 统一对所有 $n\ge2$：$a>-\dfrac12$；固定 $n$ 时：$a>-\dfrac{n-1}{2}$；不等式证明见解析
\end{itemize}
